\documentclass{article}
\usepackage{fullpage}
\usepackage{listings}
\usepackage{color}
\usepackage{amsmath,amssymb,amsthm}
\usepackage{tocloft}
\usepackage{hyperref}
\usepackage{enumitem}

\definecolor{mygreen}{rgb}{0,0.6,0}
\definecolor{mygray}{rgb}{0.5,0.5,0.5}
\definecolor{mymauve}{rgb}{0.58,0,0.82}

\lstset{
  backgroundcolor=\color{white},
  basicstyle=\ttfamily,
  breakatwhitespace=false,
  breaklines=true,
  captionpos=b,
  commentstyle=\color{mygreen},
  deletekeywords={...},
  escapeinside={\%*}{*)},
  extendedchars=true,
  frame=single,
  keepspaces=true,
  keywordstyle=\color{blue},
  morekeywords={*,...},
  numbers=left,
  numbersep=9pt,
  numberstyle=\tiny\color{mygray},
  rulecolor=\color{black},
  showspaces=false,
  showstringspaces=false,
  showtabs=false,
  stepnumber=1,
  stringstyle=\color{mymauve},
  tabsize=2,
  title=\lstname
}

\usepackage{fourier}
\usepackage{stmaryrd}

\title{
  \textbf{CSE 4392 / CSE 5369 NLP Research Project Specification} \\
}
\author{
  \normalsize
  \begin{tabular}{rcl}
    Course Instructor: & Kenny Zhu & \texttt{kenny.zhu@uta.edu} \\
    Teaching Assistant: & Peter Zhang & \texttt{wxz9630@mavs.uta.edu} \\
  \end{tabular}
}

\renewcommand{\cftsecfont}{\bfseries}
\renewcommand{\cftsecpagefont}{\bfseries}
\renewcommand{\cftsecleader}{\cftdotfill{\cftdotsep}}
\setlength{\cftbeforesecskip}{0.5em}
\setlength{\cftsecindent}{1.5em}

\newcommand{\result}[1]{\textcolor{blue}{\texttt{#1}}}
\date{Due Date: May 6, 2026, 11:59 PM}

\begin{document}

\maketitle

% \tableofcontents
% \newpage

\section{Introduction}

The course project for this class is an \emph{NLP research project}.
Each team must design and develop a system or model
that solves an interesting problem in Natural Language Processing (NLP).
The project should be research-oriented rather than purely engineering-oriented:
you are expected to formulate a clear problem, motivate its importance, implement
your own method, and evaluate it empirically. The goal is to simulate a small-scale research workflow similar to an NLP conference submission.\\

\noindent A complete project submission must include:
\begin{itemize}
  \item a working implementation,
  \item the code and data used in the project,
  \item and a research paper written in ACL format.
\end{itemize}



\section{Project Objective}

Your project must address a concrete NLP task or research question.
Examples include, but are not limited to:

\begin{itemize}
  % \item text classification,
  % \item sequence labeling,
  % \item parsing,
  % \item information extraction,
  % \item summarization,
  % \item question answering,
  % \item machine translation,
  % \item dialogue systems,
  % \item evaluation of language models,
  % \item interpretability / explainability for NLP models,
  % \item low-resource or multilingual NLP,
  % \item speech or multimodal language processing.
  \item Text Classification
  \item Named Entity Recognition (NER)
  \item Machine Translation
  \item Text Generation
  \item Question Answering 
  \item Speech Recognition
  \item Text Summarization
  \item Sentiment Analysis
  \item Information Extraction
  \item Dialogue Systems
  % \item Language Understanding (Modeling)
\end{itemize}

\section{Core Requirements}

\subsection{Original Implementation}

The project must be \textbf{your own implementation}.
You may build on public libraries and frameworks,
but the overall system, experimental pipeline, and research design must be your own work.

\begin{itemize}
  \item Copying an existing repository and only changing configuration files is not sufficient.
  \item Re-implementing a baseline from a paper is allowed, provided that you also contribute
  your own analysis, comparison, extension, or improvement.
  \item You must clearly state which components are borrowed from public toolkits
  and which parts were implemented by you.
\end{itemize}

\subsection{Research Component}

Your project must contain a genuine research component.
At minimum, your paper should answer the following questions:

\begin{itemize}
  \item What problem are you solving?
  \item Why is this problem important or interesting?
  \item What data do you need and how will you obtain it?
  \item What method or system did you design?
  \item How does it compare with one or more baselines?
  % \item What can we learn from the experimental results?
  \item How will you evaluate your method and what insights do you expect to gain?
\end{itemize}

\noindent A project that only demonstrates software usage without a research question will not receive full credit.

\subsection{Experimental Evaluation}

Every project must include an empirical evaluation.

\begin{itemize}
  \item You must evaluate your method on data that is clearly described in the paper.
  \item You must report appropriate metrics for the task.
  \item You must include at least one baseline for comparison.
  \item You should provide analysis beyond a single result table whenever possible.
\end{itemize}

\noindent Examples of further analysis include:
error analysis, ablation study, qualitative examples, comparison across data conditions,
or discussion of failure cases.

\section{Suggested Project Workflow}

A strong project will typically include the following stages:

\begin{enumerate}[label=\arabic*.]
  \item \textbf{Problem Definition} \\
  Define the NLP task clearly and precisely.

  \item \textbf{Related Work Review} \\
  Summarize prior methods and identify what gap or perspective your project addresses.

  \item \textbf{Data Preparation} \\
  Collect, clean, annotate, or preprocess the data as needed.

  \item \textbf{Method Design} \\
  Design and implement your model, algorithm, or system.

  \item \textbf{Baseline Construction} \\
  Implement or adapt one or more baseline methods for comparison.

  \item \textbf{Evaluation} \\
  Run experiments, compute metrics, and analyze the results.

  \item \textbf{Paper Writing} \\
  Present the problem, method, experiments, and findings in ACL paper format.
\end{enumerate}

\section{Paper Requirements}

\subsection{Format}

You must write your final paper in \textbf{ACL format} using the official ACL style files and templates; see the official ACL paper formatting guidelines here: \url{https://acl-org.github.io/ACLPUB/formatting.html} \\

\noindent Your paper should resemble a conference-style research paper and will typically include, but is not limited to, the following sections:

\begin{itemize}
  \item Abstract
  \item Introduction
  \item Method
  \item Experiments and Evaluation
  \item Results and Analysis
  \item Related Work
  \item Conclusion
  \item Limitations
  \item References
\end{itemize}

\subsection{Length}

The paper should target the ACL long paper style:
\begin{itemize}
  \item \textbf{8 pages of main content}, in ACL format,
  \item plus additional pages for references,
  \item and optional appendix / supplementary material if needed.
\end{itemize}

\noindent Your goal is to produce a paper of approximately conference-paper quality and scope.

\subsection{Writing Expectations}

Your paper must:

\begin{itemize}
  \item clearly describe the background, motivation, and contributions,
  \item explain the method in sufficient technical detail,
  \item describe the dataset(s) and preprocessing,
  \item report experimental settings and evaluation metrics,
  \item include quantitative results,
  \item include interpretation or analysis of the results,
  \item discuss the limitations of your work,
  \item cite prior work appropriately.
\end{itemize}

\section{Implementation Requirements}

\subsection{Programming Language and Tools}

You may use any reasonable programming language or framework for implementation,
such as Python, Java, PyTorch, TensorFlow, or other NLP toolkits.

\subsection{Reproducibility}

Your submission must be organized so that the grader can understand and reproduce your work. \\

\noindent At minimum, your code package should include:
\begin{itemize}
  \item source code,
  \item \texttt{README.md} file,
  \item instructions for the full pipeline (e.g., data preparation, training, inference, and evaluation),
  \item dependency information,
  \item and a description of where the data comes from.
\end{itemize} 

\noindent If your experiments require substantial computation, provide:
\begin{itemize}
  \item preprocessed data if allowed,
  \item saved checkpoints if helpful,
  \item and enough scripts so that the pipeline is understandable.
\end{itemize}

\subsection{Data}

You must submit the data used in the project whenever legally and practically possible.

\begin{itemize}
  \item If the dataset is public, provide download instructions and the exact version used.
  \item If the dataset is too large to submit directly, provide scripts and instructions for obtaining it.
  \item If you create your own data, clearly document how it was collected and labeled.
\end{itemize}

\section{Deliverables}

Each team should submit a single \texttt{.zip} file named \texttt{LastName\_StudentID\_Project.zip}; all team members' names and student IDs must be listed in the top-level \texttt{README.md}.\\

\noindent Your final submission must contain the following:

\begin{itemize}
  \item \textbf{Paper:} a PDF paper in ACL format, along with the LaTeX source files and any required style/template files.
  \item \textbf{Source Code:} all code needed to implement and run the project.
  \item \textbf{Data:} the dataset, processed data, or instructions/scripts to obtain the data.
  \item \textbf{README:} clear instructions for setup and execution, as well as the names and student IDs of all team members.
\end{itemize}

\noindent Before submission, make sure:
\begin{itemize}
  \item the paper compiles correctly,
  \item the code is organized,
  \item the instructions are complete,
  \item and all files are named clearly.
\end{itemize}
A good submission folder may look like this:

\begin{lstlisting}[language=bash]
LastName_StudentID_Project/
|-- paper/
|   |-- project_paper.pdf
|   `-- source_tex_files/
|-- code/
|   |-- src/
|   |-- scripts/
|   |-- train.py
|   |-- eval.py
|   |-- model.py
|   |-- requirements.txt
|   `-- ...
|-- data/
|   |-- README.md
|   `-- ...
|-- results/
|   |-- tables/
|   `-- figures/
`-- README.md
\end{lstlisting}

\section{Project Grading Criteria}

The project will be graded based on the following dimensions:

\begin{enumerate}[label=\arabic*.]
  \item \textbf{Problem Quality and Motivation} \\
  Is the problem meaningful, well-defined, and appropriate for NLP research?

  \item \textbf{Technical Depth} \\
  Does the project demonstrate real modeling, system design, or algorithmic effort?

  \item \textbf{Implementation Quality} \\
  Is the system complete, organized, and reproducible?

  \item \textbf{Experimental Design} \\
  Are the baselines, metrics, and comparisons appropriate?

  \item \textbf{Analysis and Insight} \\
  Does the paper go beyond reporting numbers and provide interpretation?

  \item \textbf{Paper Quality} \\
  Is the paper well-structured, clearly written, and formatted in ACL style?
\end{enumerate}

\section{Project Milestones}

The project will involve the following key milestones:

\begin{itemize}
  \item \textbf{Milestone 1: Proposal (Jan. 28)} \\
  Project topic, task definition, motivation, and initial related work.

  \item \textbf{Milestone 2: Progress Check (Mar. 23)} \\
  Current progress on data preparation, method development, and preliminary experiments.

  \item \textbf{Milestone 3: Final Submission (May. 6)} \\
  Final paper, code, data, and documentation.
\end{itemize}

\noindent The instructor or TA may announce concrete deadlines separately.

\section{Academic Integrity}

All submitted work must be the original work of your team.

\begin{itemize}
  \item You must properly cite papers, codebases, datasets, and external resources.
  \item Plagiarism, fabricated results, or undisclosed copying of code/text will result in a score of zero and may trigger additional academic penalties.
\end{itemize}

\section{Use of External Models and LLMs}

You may use pretrained models, APIs, or large language models as part of your project,
but you must clearly document their role. \\

\noindent For example, you should specify:
\begin{itemize}
  \item whether the model is used as a baseline, feature extractor, or core method,
  \item whether prompting, fine-tuning, or retrieval augmentation is involved,
  \item and what part of the implementation is still your own contribution.
\end{itemize}

\noindent Projects that rely entirely on black-box models with little methodological contribution
may receive limited credit unless accompanied by strong experimental analysis or a novel research question.

\section{Reference Information}

You are strongly encouraged to consult the ACL 2026 main conference submission information
for the expected paper style, length, and research presentation standard:

\begin{center}
\url{https://2026.aclweb.org/calls/main_conference_papers/#paper-submission-information}
\end{center}

\noindent This course project does not necessarily require an actual ACL submission,
but your report should be written to resemble a serious ACL-style research paper.

\end{document}